\section{Even Fibonacci Numbers --- Problem 2}
\label{prob:2}
\noindent\textbf{Source:} \url{https://github.com/KO0I/projectEuler/blob/reorganize/solutions-by-language/solutions/c/even_fibonacci_numbers_2.c}

\noindent\textbf{Assessment:} Undefined behavior.

\subsection*{Mathematical goal}
The Fibonacci recurrence is \(F_n=F_{n-1}+F_{n-2}\). The required sum includes even terms at most four million:
\[
S=\sum_{\substack{F_n\leq4\,000\,000\\2\mid F_n}}F_n.
\]

\subsection*{How the C code works}
The program allocates \verb|fib[33]|, initializes the first two entries to one and fills the array using \verb|fib[i-2]+fib[i-1]|. A second loop checks evenness with \verb|fib[k]%2==0|, accumulating qualifying entries in \verb|sum|. This separates generating the sequence from filtering it.

\subsection*{Correctness and limitations}
The declared \verb|unsigned long sum| is never initialized before \verb|sum += fib[k]|, so the behavior is undefined and the output cannot be trusted. It should start at zero. A value-bound loop would also be clearer than a fixed term count. The intended mathematical result is \texttt{...3732}, not a reliable result of this source.


\subsection*{Correction to historical C --- Problem 2}
{\color{red}\textbf{Correction:} The original accumulator was never initialized. The corrected routine begins with zero and generates Fibonacci terms until their value exceeds four million, rather than trusting a fixed array length.}
